% GENERATED FILE. Edit canonical lessons, then run npm run generate:books.
% canonical-source-hash: fnv1a64:7383e775f4a3b80f
% canonical-lessons: PA-C148-lukna, PA-C148-diggna, PA-C148-sikhauna, PA-C148-cukkna, PA-C148-tappna

\chapter{To hide, To fall, To teach, To lift, To jump}
\label{ch:pa-to-hide-to-fall-to-teach-to-lift-to-jump}
\hlchaptermodality{148}

\begin{chapteropening}
\textbf{By the end of this chapter:} \emph{I can say to hide, to fall, to teach, to lift and to jump.}

Complete the last lesson of chapter 148: I can say to hide, to fall, to teach, to lift and to jump.
\end{chapteropening}

\section[\texorpdfstring{\pa{ਲੁਕਣਾ} (lukṇā)}{ਲੁਕਣਾ (lukṇā)}]{\pa{ਲੁਕਣਾ} (lukṇā) — to hide}

\label{lesson:PA-C148-lukna}



\begin{quote}
Before the new one: say the Punjabi for strong, then the Punjabi for weak.

Read the form line once: \textbf{\pa{ਕੰਮ}: \pa{ਨੌਕਰੀ}}. What does this person do? (A job, \emph{naukarī}.)
\end{quote}

\subsection*{You'll want to know: \pa{ਲੁਕਣਾ}}
\textbf{\pa{ਲੁਕਣਾ}} — \emph{lukṇā} — \textquotedblleft{}to hide\textquotedblright{}.

\begin{grammarlens}[title={What you've built}]
One more word for this chapter.
\end{grammarlens}

\subsection*{Your turn}
\begin{itemize}
\raggedright
  \item \emph{Say it:} \emph{lukṇā}
  \item \emph{Say it:} \emph{lukṇā}, once more
  \item \emph{Say it:} say \emph{lukṇā}
  \item \emph{Recall:} say the Punjabi for strong, then the Punjabi for weak, then say \emph{lukṇā} again
\end{itemize}

\begin{tcolorbox}[breakable,colback=teal!4,colframe=teal!35!black,title={Before you move on}]
What does \pa{ਲੁਕਣਾ} mean? (\textquotedblleft{}To hide\textquotedblright{}.) Say it once more.
\end{tcolorbox}
\section[\texorpdfstring{\pa{ਡਿੱਗਣਾ} (ḍiggṇā)}{ਡਿੱਗਣਾ (ḍiggṇā)}]{\pa{ਡਿੱਗਣਾ} (ḍiggṇā) — to fall}

\label{lesson:PA-C148-diggna}



\begin{quote}
Before the new one: say the Punjabi for weak, then the Punjabi for to hide.
\end{quote}

\subsection*{You'll want to know: \pa{ਡਿੱਗਣਾ}}
\textbf{\pa{ਡਿੱਗਣਾ}} — \emph{ḍiggṇā} — \textquotedblleft{}to fall\textquotedblright{}.

\begin{grammarlens}[title={What you've built}]
One more word for this chapter.
\end{grammarlens}

\subsection*{Your turn}
\begin{itemize}
\raggedright
  \item \emph{Say it:} \emph{ḍiggṇā}
  \item \emph{Say it:} \emph{ḍiggṇā}, once more
  \item \emph{Say it:} say \emph{ḍiggṇā}
  \item \emph{Recall:} say the Punjabi for weak, then the Punjabi for to hide, then say \emph{ḍiggṇā} again
\end{itemize}

\begin{tcolorbox}[breakable,colback=teal!4,colframe=teal!35!black,title={Before you move on}]
What does \pa{ਡਿੱਗਣਾ} mean? (\textquotedblleft{}To fall\textquotedblright{}.) Say it once more.
\end{tcolorbox}
\section[\texorpdfstring{\pa{ਸਿਖਾਉਣਾ} (sikhāuṇā)}{ਸਿਖਾਉਣਾ (sikhāuṇā)}]{\pa{ਸਿਖਾਉਣਾ} (sikhāuṇā) — to teach}

\label{lesson:PA-C148-sikhauna}



\begin{quote}
Before the new one: say the Punjabi for to hide, then the Punjabi for to fall.
\end{quote}

\subsection*{You'll want to know: \pa{ਸਿਖਾਉਣਾ}}
\textbf{\pa{ਸਿਖਾਉਣਾ}} — \emph{sikhāuṇā} — \textquotedblleft{}to teach\textquotedblright{}.

\begin{grammarlens}[title={What you've built}]
One more word for this chapter.
\end{grammarlens}

\subsection*{Your turn}
\begin{itemize}
\raggedright
  \item \emph{Say it:} \emph{sikhāuṇā}
  \item \emph{Say it:} \emph{sikhāuṇā}, once more
  \item \emph{Say it:} say \emph{sikhāuṇā}
  \item \emph{Recall:} say the Punjabi for to hide, then the Punjabi for to fall, then say \emph{sikhāuṇā} again
\end{itemize}

\begin{tcolorbox}[breakable,colback=teal!4,colframe=teal!35!black,title={Before you move on}]
What does \pa{ਸਿਖਾਉਣਾ} mean? (\textquotedblleft{}To teach\textquotedblright{}.) Say it once more.
\end{tcolorbox}
\section[\texorpdfstring{\pa{ਚੁੱਕਣਾ} (cukkṇā)}{ਚੁੱਕਣਾ (cukkṇā)}]{\pa{ਚੁੱਕਣਾ} (cukkṇā) — to lift}

\label{lesson:PA-C148-cukkna}



\begin{quote}
Before the new one: say the Punjabi for to fall, then the Punjabi for to teach.
\end{quote}

\subsection*{You'll want to know: \pa{ਚੁੱਕਣਾ}}
\textbf{\pa{ਚੁੱਕਣਾ}} — \emph{cukkṇā} — \textquotedblleft{}to lift\textquotedblright{}.

\begin{grammarlens}[title={What you've built}]
One more word for this chapter.
\end{grammarlens}

\subsection*{Your turn}
\begin{itemize}
\raggedright
  \item \emph{Say it:} \emph{cukkṇā}
  \item \emph{Say it:} \emph{cukkṇā}, once more
  \item \emph{Say it:} say \emph{cukkṇā}
  \item \emph{Recall:} say the Punjabi for to fall, then the Punjabi for to teach, then say \emph{cukkṇā} again
\end{itemize}

\begin{tcolorbox}[breakable,colback=teal!4,colframe=teal!35!black,title={Before you move on}]
What does \pa{ਚੁੱਕਣਾ} mean? (\textquotedblleft{}To lift\textquotedblright{}.) Say it once more.
\end{tcolorbox}
\section[\texorpdfstring{\pa{ਟੱਪਣਾ} (ṭappṇā)}{ਟੱਪਣਾ (ṭappṇā)}]{\pa{ਟੱਪਣਾ} (ṭappṇā) — to jump}

\label{lesson:PA-C148-tappna}



\begin{quote}
Before the new one: say the Punjabi for to teach, then the Punjabi for to lift.
\end{quote}

\subsection*{You'll want to know: \pa{ਟੱਪਣਾ}}
\textbf{\pa{ਟੱਪਣਾ}} — \emph{ṭappṇā} — \textquotedblleft{}to jump\textquotedblright{}.

\begin{grammarlens}[title={What you've built}]
One more word for this chapter.
\end{grammarlens}

\subsection*{Your turn}
\begin{itemize}
\raggedright
  \item \emph{Say it:} \emph{ṭappṇā}
  \item \emph{Say it:} \emph{ṭappṇā}, once more
  \item \emph{Say it:} say \emph{ṭappṇā}
  \item \emph{Recall:} say the Punjabi for to teach, then the Punjabi for to lift, then say \emph{ṭappṇā} again
\end{itemize}

\begin{tcolorbox}[breakable,colback=teal!4,colframe=teal!35!black,title={Before you move on}]
What does \pa{ਟੱਪਣਾ} mean? (\textquotedblleft{}To jump\textquotedblright{}.) Say it once more.
\end{tcolorbox}
